\section{引言}
SiZrOC 陶瓷纤维因其优异的高温稳定性与抗氧化性能，在航空航天及耐高温复合材料领域具有广泛的应用前景。本实践项目主要围绕 SiZrOC 纤维的表面改性展开探索。针对单一成分纤维在结构与力学性能上的局限性，本研究尝试引入多层静电纺丝技术，旨在通过控制不同层级的物理化学性质，优化纤维的前驱体成形质量及烧结后的界面结合特性。

\section{实验设计与制备}
在多层纺丝工艺的探索中，本实验拟定了两种实现方案：微观层面的同轴套管样式与宏观层面的"夹心饼干"叠层样式。受限于同轴挤出针头的硬件条件，本阶段主要采用宏观三层叠层方案进行实验。

实验采用有机层与无机层交替的方法进行表层成形。为维持有机网络的均匀性，在有机前驱体中加入了 1.72\,g 至 1.735\,g 的 PVP 作为纺丝助剂。在优化后的三层制备方案中，首底两层均采用含 1.735\,g PVP 的材料，而中间层则额外添加了 2\,g 葡萄糖与 0.64\,g 正丙醇锆。这种组分调节使得中间层表现出较强的粘性与一定的韧性。前驱体成形后，按照预设的程序（200/400/420/1250/60/−121）进行高温烧结处理。

\section{结果与讨论}
在多层前驱体的揭样剥离过程中，通过镊子可以轻易地将三层结构分离，且未观察到明显的粘滞拖拽现象，这表明各层之间保持了良好的独立性，层间界面清晰。烧结后产物的宏观形貌与前期单层样品无显著差异。

为了进一步观察纤维的微观形貌与表面改性效果，对样品进行了 SEM（扫描电子显微镜）检测。然而初步测试结果并不理想，经分析认为，测试环境（107 实验室）湿度偏高是导致形貌缺陷或成纤不良的主要原因。随后，实验设备被整体搬迁至环境控制更佳的创新大楼 A 座 201 房间，并重新完善了静电纺丝机及高压电源的接地等电气保障措施，为后续的 5321 硅烷原液改性实验（配比：15+15+0.96）奠定了基础。

\section{结论}
本研究初步验证了基于"夹心饼干"模式的三层静电纺丝工艺在 SiZrOC 纤维制备中的可行性。通过添加 PVP、葡萄糖与正丙醇锆，成功调节了中间层的粘韧性，实现了界面清晰的多层前驱体制备。后续实验将在优化后的低湿度环境中，进一步探究硅烷偶联剂对纤维界面的深度改性作用。
